\chapter{Lab 1 -- Introduzione a PyTorch con una FFNN}
\label{cap:lab1}

\begin{nota}[Fonti e limiti di questo capitolo]
Testo del laboratorio ``\emph{1 -- Introduction to Deep Learning: building, training, and judging a
Feed-Forward Neural Network}'' pubblicato sul portale. La consegna è \textbf{parafrasata}; per i
dettagli fa fede il testo originale. Il Lab~1 non dà voto ma è obbligatorio e vale il bonus: come
per gli altri lab qui \textbf{non} ci sono soluzioni, codice, numeri o scelte di un gruppo.
\end{nota}

\section{Contesto e obiettivi}
Il lab ha due scopi: introdurre PyTorch e il comportamento delle reti neurali su un vero problema
di classificazione, e dare una prima esperienza (non valutata) di lab e report come allenamento per
quelli valutati. Il testo insiste che è un lab di \textbf{giudizio, non di ricette}: si riceve un
obiettivo, dei vincoli e un punto di partenza consigliato, e si devono decidere, giustificare e
difendere le scelte sulla base di ciò che si osserva. In concreto si vuole:
\begin{itemize}
  \item costruire e addestrare FFNN per \en{intrusion detection} con PyTorch;
  \item vedere come il training si rompe quando un fondamentale (scaling, learning rate) è sbagliato
        e imparare a leggerne i sintomi;
  \item gestire lo sbilanciamento e ragionare su quale metrica conta per un IDS;
  \item vedere come una singola feature può guidare di nascosto il modello e come una scelta di
        pulizia può rimodellare le classi;
  \item costruire una baseline semplice e verificare se la complessità aggiuntiva ne vale la pena.
\end{itemize}

\section{Dataset}
Un sottoinsieme di \textbf{CICIDS2017} (Canadian Institute for Cybersecurity, University of New
Brunswick), benchmark di intrusion detection con traffico benigno realistico e attacchi comuni. Un
unico CSV di flussi etichettati; ogni riga è un flusso descritto da poche feature a livello di flusso
(durate e tempi di interarrivo, statistiche sulle lunghezze dei pacchetti, rate di byte e pacchetti,
alcuni contatori di flag TCP, numero di pacchetti in avanti, porta di destinazione, \ldots). La colonna
\code{Label} ha quattro valori: \code{Benign}, \code{PortScan}, \code{DoS Hulk}, \code{Brute Force}.
Il testo \textbf{non} dice di proposito quanto sono bilanciate le classi né quali problemi abbiano i
dati: scoprirlo è il primo task.

\section{Consegna (parafrasi)}
Regole di lavoro comuni a tutto il lab: logbook con una riga per run, budget di circa 25 run,
ipotesi prima di ogni run, una modifica alla volta (cap.~\ref{cap:lab-report}).

\begin{description}[style=nextline,leftmargin=1.2em]
\item[Task 1 -- Preprocessing (nessun run)]
  Ispezionare i dati (numero di campioni, classi e bilanciamento, problemi di qualità come valori
  mancanti, duplicati, valori implausibili); decidere cosa fare di mancanti e duplicati e dire quanti
  campioni restano; scegliere le proporzioni dello split train/validation/test e se stratificare;
  scegliere e motivare lo scaling dalle distribuzioni osservate (la porta di destinazione, per ora,
  va trattata come numerica). Domande: perché stratificare o no, su quale partizione si fitta lo
  scaler e cosa va storto se lo si fitta su tutto, come gli outlier guidano la scelta.
\item[Task 2A -- Un primo classificatore]
  Una FFNN a \textbf{un solo hidden layer} per le 4 classi, con un training loop pulito (loss,
  ottimizzatore, early stopping sulla validation). Il testo suggerisce un punto di partenza semplice
  (un hidden layer piccolo con ReLU, cross-entropy, AdamW, learning rate e batch size moderati, early
  stopping) e chiede di dichiarare ogni scelta. Prima di addestrare: prevedere quali classi saranno
  più difficili e perché. Poi curve di loss e classification report di validation; come si sceglie il
  modello ``migliore'' tra le epoch, se le curve convergono, se le previsioni reggono.
\item[Task 2B -- Perché i fondamentali contano]
  Tre varianti della configurazione 2A, una sola cosa sbagliata alla volta: S1 senza scaling, S2 con
  learning rate troppo alto (fino a rompere il training), S3 con learning rate molto più basso. Per
  ognuna: prevedere la curva, plottare train/validation loss, dare un nome al \en{failure mode} e
  spiegarne il meccanismo. Poi ordinarle per gravità, dire quale fondamentale sistemare per primo su
  un problema nuovo e cosa, nella sola forma di una curva, indicherebbe la causa.
\item[Task 2C -- Fissare la baseline]
  Scegliere un modello a un hidden layer ben tarato e compreso: è la ``baseline da battere''.
  Classification report di validation e test, commento sul gap. Come esperimento controllato, cambiare
  l'attivazione (es.\ lineare $\leftrightarrow$ ReLU) e discutere se sia lecito confrontare le due
  varianti sul test set.
\item[Task 3 -- Una feature che distorce il modello]
  Nei dati, tutti i flussi \code{Brute Force} hanno porta di destinazione 80. Domanda: è una proprietà
  degli attacchi brute force o un artefatto del dataset? Esperimento controfattuale \textbf{solo sul
  test}: cambiare la porta di quei flussi, prevedere l'effetto sulla recall e verificarlo con la
  baseline senza riaddestrare. Poi rimuovere la porta dai dati originali, rifare il preprocessing e
  usare il dataset senza porta per il resto del lab. Domande sul conteggio delle classi dopo la
  pulizia, su quale classe si riduce e perché, su quale diventa la più rara e sul dilemma: de-duplicare
  ancora o no?
\item[Task 4 -- La classe rara: loss e sbilanciamento]
  Riaddestrare la baseline sui dati senza porta, prevedere quale classe soffre di più, poi usare una
  cross-entropy \textbf{pesata per classe}. Domande: da quale partizione si calcolano i pesi e perché
  ogni altra scelta è sbagliata; accuracy e macro-F1 si muovono nella stessa direzione? Quale metrica
  ottimizzare per un IDS che deve trovare attacchi rari?
\item[Task 5 -- Iperparametri (circa 10 run in tutto)]
  5A profondità (2--3 candidati più profondi scelti e motivati, niente grid); 5B batch size (tre
  valori che ``circondano'' un'ipotesi, effetto su generalizzazione e tempo); 5C ottimizzatore contro
  learning rate (SGD, SGD+momentum, AdamW variando il learning rate: chi domina?); 5D verdetto sul test
  del miglior modello affiancato alla baseline: la complessità ne valeva i run?
\item[Task 6 -- Ottimizzazioni]
  Una rete grande e profonda imposta dal testo, addestrata a lungo senza regolarizzazione né early
  stopping: prevedere le curve e mostrare dove si vede l'overfitting. Poi, una modifica per run e dalla
  più semplice: rete più piccola o stop anticipato, poi un regolarizzatore alla volta (weight decay,
  dropout, batch normalization), controllando la validation prima di aggiungere il successivo.
  Domande: cosa ha aiutato di più con meno sforzo; la rete grande regolarizzata batte la baseline
  semplice sul test?
\item[Riflessione finale]
  Guardando tutto il logbook: quali decisioni hanno dato più prestazioni per unità di sforzo?
\end{description}

\section{Teoria coinvolta}
\begin{itemize}
  \item Pulizia, duplicati, valori mancanti, split stratificato, scaling e leakage; struttura di una
        FFNN e training loop in PyTorch: cap.~\ref{cap:ffnn}.
  \item Loss, cross-entropy, ottimizzatori (SGD, momentum, Adam/AdamW), learning rate, batch size,
        metriche per classe, learning curve, under/overfitting: cap.~\ref{cap:addestramento}.
  \item Attivazioni e collasso lineare di una rete senza non linearità: cap.~\ref{cap:reti-neurali}.
  \item Early stopping, weight decay, dropout, batch normalization: cap.~\ref{cap:regolarizzazione}.
  \item Feature-scorciatoia (\en{shortcut}) e porte come feature categoriche: cap.~\ref{cap:ffnn}
        e~\ref{cap:embedding}.
\end{itemize}

\section{Metodo generale}
\begin{enumerate}
  \item \textbf{Prima i dati}: contare, guardare distribuzioni e code, trovare duplicati e anomalie di
        qualità; ogni passo di pulizia va motivato e il numero di righe prima/dopo va riportato.
  \item \textbf{Split prima di tutto il resto}; fit di scaler e pesi di classe sul solo training.
  \item \textbf{Baseline semplice} che si capisce fino in fondo, con curve e report per classe.
  \item \textbf{Una variabile alla volta}, con ipotesi scritta prima: ogni esperimento risponde a una
        domanda precisa del testo.
  \item \textbf{Validation per scegliere, test per verificare}: il test non entra mai nelle scelte.
  \item \textbf{Confronto onesto} con la baseline: la complessità si giustifica solo con un
        miglioramento che regge e che vale il costo.
\end{enumerate}

\section{Insidie tipiche}
\begin{attenzione}[errori da evitare]
\begin{itemize}
  \item Fittare lo scaler (o calcolare i pesi di classe) su tutto il dataset o sul validation/test.
  \item Giudicare dalla sola accuracy un problema sbilanciato: guardare report per classe e macro-F1.
  \item Scegliere tra varianti (attivazione, architettura) guardando il test set.
  \item Riportare il modello dell'ultima epoch invece di quello scelto dall'early stopping.
  \item Cambiare due cose nello stesso run, o fare grid search ``a tappeto'' bruciando il budget.
  \item Spiegare i risultati con frasi da manuale invece che con le proprie curve e i propri numeri.
  \item Dimenticare che rimuovere una feature e de-duplicare possono cambiare la composizione delle
        classi: ricontare dopo ogni modifica.
  \item Curve illeggibili o senza legenda; affermazioni ``è meglio'' senza metrica di riferimento.
\end{itemize}
\end{attenzione}

\section{Checklist}
\begin{itemize}
  \item[$\square$] Numero di campioni prima e dopo la pulizia, con motivazione di ogni scelta.
  \item[$\square$] Split e scaling motivati; dove è stato fittato lo scaler.
  \item[$\square$] Previsioni scritte prima di ogni run e confrontate dopo.
  \item[$\square$] Le tre curve di 2B con nome del failure mode e meccanismo.
  \item[$\square$] Discussione validation vs test della baseline; perché il test non serve a scegliere.
  \item[$\square$] Esperimento controfattuale sulla porta e ricalcolo delle classi senza porta.
  \item[$\square$] Partizione da cui si calcolano i pesi di classe; scelta motivata della metrica.
  \item[$\square$] Verdetto onesto su profondità, batch size, ottimizzatori e regolarizzazione rispetto
        alla baseline.
  \item[$\square$] Logbook entro il budget; riflessione finale argomentata dai propri run.
  \item[$\square$] Checklist generale del cap.~\ref{cap:lab-report}.
\end{itemize}
